\chapter{Algebraic Branches and Rational Curves}\label{ch:algebra-fta}

An algebraic formula is useful only after a branch has been selected and
its values enclosed.  We develop representative constructions: positive
square roots, a certified simple complex root, and a rational curve.  Each
has its own finite domain and separation data.

\section{Positive roots and their derivatives}
On $[a,b]$ with rational $a>0$, the square-root bisection computation has
\[
 \frac{\sqrt y-\sqrt x}{y-x}
      =\frac1{\sqrt y+\sqrt x}\qquad(x\ne y).
\]
As both inputs increase these slopes decrease, so the function is concave.
On the interval the denominator is at least $2\sqrt a$, with an explicit
rational positive lower bound obtainable by refinement.  The difference
between the left and right secants at $x$ tends to zero by the root
continuity estimate.  Consequently
\[
 D\sqrt x=\frac1{2\sqrt x}.
\]
The derivative is not bounded near zero; no theorem requiring a bounded
endpoint derivative is applied there.

For a primitive of the square root, set $F(x)=\frac23 x\sqrt x$ for
$x\ge0$, and extend by zero for $x\le0$.  If $0\le x<y$, write
$u=\sqrt x$, $v=\sqrt y$.  Then
\[
 \frac{F(y)-F(x)}{y-x}
       =\frac{2(v^2+uv+u^2)}{3(u+v)}
\]
when $u+v>0$.  It lies between $u$ and $v$: after clearing the positive
denominator the two differences are
$(v-u)(2v+u)$ and $(v-u)(v+2u)$.  Adjacent secants are therefore ordered;
$F$ is convex and $DF(x)=\sqrt x$ on $x\ge0$, including the shrinking
one-sided bound at zero.  The convex FTC gives
\[
 \int_0^b\sqrt x\,dx=\frac23 b\sqrt b.
\]
The reciprocal-root integral can also be computed as an improper integral:
removing $[0,\varepsilon]$ has error $2\sqrt\varepsilon$, by the
primitive $2\sqrt x$ on positive interior intervals and their FTC.  The tail bound, not a bare
improper-integral symbol, defines the computation.

\section{Separated rational functions have finite remainder bounds}
Let $P,Q$ be rational polynomials on a fixed interval with the supplied
separation $|Q|\ge c>0$.  For numbers $v,w$ separated from zero,
\[
 \frac1w-\frac1v+\frac{w-v}{v^2}
       =\frac{(w-v)^2}{v^2w}.
\]
Finite polynomial expansions and this identity give
\[
 \frac{P(x+h)}{Q(x+h)}-\frac{P(x)}{Q(x)}
 =\frac{P'(x)Q(x)-P(x)Q'(x)}{Q(x)^2}h+R(x,h),
 \qquad |R(x,h)|\le K h^2,
\]
with a rational $K$ computed from polynomial coefficient bounds, $c$, and
the interval length.  Indeed the reciprocal remainder is bounded by
$c^{-3}|Q(x+h)-Q(x)|^2$, and the polynomial linear remainders are bounded
by their finite binomial expansions.  Multiplication adds only bounded
quadratic terms.

Applying the estimate at $h$ and $-h$ cancels the linear terms.  Thus
$P/Q+Kx^2$ is midpoint convex, hence convex, and its centered secant gaps
are at most $4K|h|$.  The quotient formula is consequently the actual
computed derivative, with a convex-difference FTC certificate.  This
supplies, in particular, the rational substitutions in this chapter and
the final elliptic-integral example.  No assertion about an arbitrary
pointwise rational second derivative is used.

\section{A complex root certified by contraction}
Let $P$ be a polynomial with rational complex coefficients, and let $c$
be a rational complex point with $P'(c)\ne0$.  Define
\[
 T(z)=z-P(z)/P'(c).
\]
Use the coordinate maximum norm on complex numbers.  Multiplication by
$w$ has norm at most $|\Re w|+|\Im w|$.  Suppose finite polynomial
bounds on the square $\|z-c\|_\infty\le r$ certify
\[
 \|T(c)-c\|_\infty\le(1-q)r,\qquad
 \|T(z)-T(w)\|_\infty\le q\|z-w\|_\infty,
 \quad 0\le q<1.
\]
The second inequality can be checked by factoring $P(z)-P(w)$ and
bounding the resulting polynomial on two copies of the square.
Then $T$ preserves the square.  Starting at $z_0=c$, compute the rational
iterates $z_{n+1}=T(z_n)$.  For $m>n$,
\[
 \|z_m-z_n\|_\infty
       \le\frac{q^n}{1-q}\|z_1-z_0\|_\infty.
\]
Rational boxes with this error, stabilized coordinatewise, define a root
computation.  Polynomial continuity and the iteration identity prove
$P(z)=0$.  Two root computations in the square have distance at most
$q^n$ times their initial distance for every $n$, proving uniqueness
without a fixed-point existence theorem.

\begin{example}[A nonrational complex root]
For $P(z)=z^2+2$, take $c=3i/2$ and $r=1/6$.  The residual step has norm
$1/12$.  Factoring $T(z)-T(w)$ gives contraction constant $q=2/9$ on
the square, since
\[
 T(z)-T(w)=(z-w)\left(1-\frac{z+w}{3i}\right).
\]
Also $1/12\le(1-2/9)/6$.  Thus the iteration constructs its unique root
there.  Comparison with the square-root construction identifies it as
$i\sqrt2$; reflection gives the other root.  The exact polynomial identity
$(z-i\sqrt2)(z+i\sqrt2)=z^2+2$ accounts for both roots.
\end{example}

This is a certificate for supplied root locations, not an invocation of a
general complex root-existence theorem.  Polynomial identities, separation,
and refinement are the concrete mathematical ingredients.

\section{Rational parametrization changes the computation}
The singular cubic
\[
 y^2=x^2(1-x)
\]
has the parametrization
\[
 x=1-t^2,\qquad y=t(1-t^2).
\]
Substitution verifies the equation by finite algebra.  Away from $x=0$,
the inverse parameter is $t=y/x$.  Both coordinates are rational at
rational $t$.  A rational differential form in $x,y$ therefore becomes a
rational function of $t$ wherever the denominators are separated from
zero.

For example, on $0<t<1$,
\[
 \frac{dx}{y}=-\frac{2\,dt}{1-t^2}.
\]
Between rational parameters $0<u<v<1$, its integral is computed by
\[
 -\log\frac{1+v}{1-v}+\log\frac{1+u}{1-u}.
\]
Indeed $-2/(1-t^2)=-1/(1-t)-1/(1+t)$, and the logarithm FTC proves the
formula.  The line-integral computation is defined along the parametrized
path, so its orientation and domain are explicit.

A parametrization is more than a closed formula: it can replace an
algebraic-value algorithm by exact rational evaluation and turn a difficult
integral into a previously understood one.  The final chapter will use a
rational circle chart to regularize an integral whose next useful description is a differential equation
rather than an elementary primitive.
